% !TEX program = lualatex
%%%%%%%%%%%%%%%%%%%%%%%%%%%%%%%%%%%%%%%%%%%%%%%%%%%%%%%%%%%%%%%%%%%%%%%%%%%%%%%
%% APRESENTAÇÃO FOCADA / DEMO DAY -- HANARO . Setembro/2026 . VERSAO 2
%% IFAM . INOVA  |  AX Academy  |  Cliente: LG
%%
%% Compilar:  uv run build demoday-focado
%%
%% ESTRUTURA -- 7 slides focados
%%   1 Capa . 2 Equipe . 3 Contextualizacao . 4 Objetivo geral .
%%   5 Sistema em Deploy (Intranet LG :4203) .
%%   6 Resultados esperados . 7 Encerramento
%%
%% NOTA: Os slides "Escopo do MVP" e "A Marca Hanaro" foram suprimidos
%% nesta versão focada para uma apresentação executiva mais direta e objetiva.
%%%%%%%%%%%%%%%%%%%%%%%%%%%%%%%%%%%%%%%%%%%%%%%%%%%%%%%%%%%%%%%%%%%%%%%%%%%%%%%

\documentclass[rodape]{hanarokickoff}

\dxiprojeto{Hanaro · Controle de Perdas (Scrap) · Apresentação Executiva}

%  A classe traz "Powered by AX Academy · Digital Transformation" na capa;
%  aqui a assinatura fica so na marca.
\dxipowered{\textit{Powered by} \textbf{AX Academy}}

%% Cartao de contribuicao do slide de resultados. Tres deles formam a fileira:
%% 3 x 8,5 + 2 x 0,9 = 27,3 cm, dentro dos 27,4 uteis.
%% \ddcontrib{<numero>}{<titulo>}{<texto>}
\newcommand{\ddcontrib}[3]{%
  \begin{minipage}[t]{8.5cm}\vspace{0pt}\RaggedRight\hyphenpenalty=10000
    {\dximont\bfseries\fontsize{26}{30}\selectfont\color{dxicinza}#1}\\[0.10em]
    {\dximont\bfseries\fontsize{15}{19}\selectfont\color{lgred}#2}\\[0.30em]
    {\dxiurb\fontsize{12.5}{17}\selectfont\color{dxinavy!85}#3}%
  \end{minipage}\hspace{0.9cm}\ignorespaces}

\begin{document}

%% ---------------------------------------------------------------- SLIDE  1 --
\dxicapa
  {APRESENTAÇÃO EXECUTIVA · SETEMBRO DE 2026}
  {\includegraphics[width=14.2cm]{marca/hanaro-assinatura-branca.pdf}}
  {Controle de Perdas (Scrap) da Linha Final: da planilha manual ao
   acompanhamento centralizado de IF Cost.}


%% ---------------------------------------------------------------- SLIDE  2 --
%% Uma linha por pessoa, nao um slide por pessoa: em 5 minutos, quatro slides de
%% biografia consomem 2 min antes de o problema ser enunciado. O texto longo vai
%% para o e-book.
\dxiequipe[Equipe do projeto]{%
  \dximembro[fotos/card-andre]{Líder Técnico · P.O.}{André Filipe Aloise}
  \dximembro[fotos/card-messyas]{Desenvolvedor FullStack}{Messyas Gois França}
  \dximembro[fotos/card-gustavo]{Eng. de Software · CI/CD}{Gustavo Almada}
  \dximembro[fotos/card-teodorio]{Desenvolvedor Frontend}{Teodório Neto}
}


%% ---------------------------------------------------------------- SLIDE  3 --
%% Tudo em topico. O destaque em vermelho e' o unico numero do slide.
\begin{dxiconteudo}{CONTEXTUALIZAÇÃO}

\vspace{-0.35cm}

\dxicol{14.6}{%
  {\dximont\bfseries\fontsize{16}{21}\selectfont\color{lgred}Situação atual}

  \vspace{0.45em}
  {\dxiCorpoP\color{dxinavy}
  \begin{dxilista}[dxinavy]
    \item Acompanhamento de \textbf{Scrap} conduzido \textbf{fora de qualquer
          sistema}.
    \item A Qualidade extrai os dados, aplica a conversão cambial e encaminha
          a planilha por correio eletrônico.
    \item A Produção reprocessa os dados manualmente em Excel antes da
          decisão.
  \end{dxilista}\par}

  \vspace{1.0em}
  {\dximont\bfseries\fontsize{15}{20}\selectfont\color{dxired}%
   1,5 a 2 h por dia, até 504 h por ano\par}
  \vspace{0.2em}
  {\dxiCorpoP\color{dxinavy!85}
   2 h/dia $\times$ 21 dias úteis $\times$ 12 meses, equivalente à ocupação
   integral de um analista em consolidação de planilha.\par}
}%
\dxicol{13.2}{%
  {\dximont\bfseries\fontsize{16}{21}\selectfont\color{lgred}Impactos identificados}

  \vspace{0.45em}
  {\dxiCorpoP\color{dxinavy}
  \begin{dxilista}[dxired]
    \item Cerca de um dia de trabalho por semana de um analista consumido em
          consolidação manual.
    \item \textbf{Inconsistência recorrente} no número reportado, sem
          visualização executiva automatizada.
    \item \textbf{VS} (MyLink) e \textbf{BM} (Módulos) permanecem fora do
          cálculo: o tempo disponível não permite apurar todas as fábricas.
  \end{dxilista}\par}
}

\end{dxiconteudo}


%% ---------------------------------------------------------------- SLIDE  4 --
\dxidestaque
  {OBJETIVO GERAL}
  {\textbf{Centralizar e monitorar} o controle de perdas}
  {Automatizar a ingestão, o tratamento e a análise dos dados de Scrap,
   entregando \textbf{uma base única e confiável} e um \textbf{painel
   executivo em tempo real} que sustente a decisão gerencial sobre o
   \textbf{IF Cost}.}


%% ---------------------------------------------------------------- SLIDE  5 --
%% Sistema em deploy na intranet LG (:4203) e validacao de PoC
\begin{dxiconteudo}[semdeco]{SISTEMA EM DEPLOY}

\vspace{-0.45cm}
\dxicol{14.6}{%
  \includegraphics[width=14.2cm]{prototipo/portal-desktop}
}%
\dxicol{12.4}{%
  {\dximont\bfseries\fontsize{15}{20}\selectfont\color{lgred}%
   Ambiente em Produção (Intranet LG)\par}

  \vspace{0.35em}
  {\dxiurb\fontsize{11.5}{15}\selectfont\color{dxinavy!85}
   Sistema implantado e acessível na rede interna:\\
   \textbf{\color{lgred}Intranet LG \textbullet\ Porta 4203}\par}

  \vspace{0.45em}
  {\dximont\bfseries\fontsize{12}{16}\selectfont\color{dxinavy}%
   Funcionalidades ativas\par}
  \vspace{0.2em}
  {\dxiurb\fontsize{11}{14.5}\selectfont\color{dxinavy!85}
  \begin{dxilista}[hanaroverde]
    \item IF Cost acumulado contra o \textbf{target}, com o \textit{gap} em
          dinheiro e em pontos percentuais.
    \item \textbf{Target $\times$ realizado} mensal, com histórico comparativo
          entre anos na mesma escala.
    \item Distribuição por \textbf{produto e área}, Top 10 de linhas, modelos
          e componentes e \textbf{Top 5 ofensores}.
    \item \textit{Toggles} de \textbf{impacto absoluto} e de \textbf{valor
          monetário}, prontos para exibição em TVs de fábrica.
  \end{dxilista}\par}

  \vspace{0.40em}
  {\dximont\bfseries\fontsize{12}{16}\selectfont\color{dxinavy}%
   Validação \& Status dos Agentes\par}
  \vspace{0.20em}
  {\dxiurb\fontsize{10}{13.5}\selectfont\color{dxinavy!85}
   $\bullet$~\textbf{Assistente ReAct (LangGraph):} Fluxo \textbf{aprovado visualmente}
   pelo Gerente da Produção (Leandro Corazza) em fase de PoC. O agente em si está em
   homologação para deploy produtivo.\\[2pt]
   $\bullet$~\textbf{Impedimento RPA:} Contas aguardam autorização de acesso ao GERP;
   ingestão operando por carga de planilha.\par}
}

\end{dxiconteudo}


%% ---------------------------------------------------------------- SLIDE  6 --
%% Continua "esperados" de proposito: o sistema ainda nao esta em producao e a
%% frente de RPA esta bloqueada, entao anunciar resultado realizado seria falso.
%% UM destaque so -- a reducao de horas -- e tres contribuicoes.
\begin{dxiconteudo}{RESULTADOS ESPERADOS}

\vspace{-0.15cm}

\dxicol{13.0}{%
  {\dximont\bfseries\fontsize{70}{76}\selectfont\color{dxired}504 h}\\[0.35em]
  {\dximont\bfseries\fontsize{17}{22}\selectfont\color{dxinavy}%
   de trabalho manual devolvidas à equipe por ano\par}
}%
\dxicol{13.4}{%
  \vspace{1.5cm}
  {\dxiurb\fontsize{14}{19}\selectfont\color{dxinavy!80}
   2 h/dia $\times$ 21 dias úteis $\times$ 12 meses, equivalente à ocupação
   integral de um analista, hoje dedicada à consolidação de planilha.\par}
}

\vspace{1.2em}
{\color{dxicinza}\rule{27.4cm}{0.6pt}}
\vspace{1.2em}

%  ARMADILHA: o rodape da classe comeca por volta de 17,4 cm. Com cinco linhas
%  de texto no cartao, a terceira coluna passa por cima da assinatura AX
%  Academy sem gerar Overfull. Quatro linhas e' o teto -- medido no PDF.
\ddcontrib{01}{Base única para o IF Cost}{%
  Elimina a divergência entre planilhas: a base consolidada passa a ser a
  fonte oficial, refletida no último \textit{snapshot} validado.}%
\ddcontrib{02}{Identificação do item ofensor}{%
  Pareto por quantidade e por IF Cost com o item identificado em lugar do
  \textit{part number}, o que qualifica a análise de causa.}%
\ddcontrib{03}{Registro da justificativa pelo gestor}{%
  O alerta é encaminhado diretamente ao gestor responsável e a resposta é
  registrada no sistema, sem intervenção do analista.}

\end{dxiconteudo}


%% ---------------------------------------------------------------- SLIDE  7 --
\dxidivisor[Cliente: LG Electronics · Gestão de Produção]
  {\fontsize{38}{46}\selectfont OBRIGADO / THANK YOU /
   {\dxihangul\bfseries 감사합니다}}
  {Hanaro · Controle de Perdas (Scrap) da Linha Final.\\
   IFAM · INOVA · AX Academy.}

\end{document}
